\documentclass[aspectratio=169]{beamer}

% =========================================================
% Кодировка и язык
% =========================================================
\usepackage{fontspec}
\usepackage[russian]{babel}
\babelfont{rm}{Liberation Serif}
\babelfont{sf}{Liberation Sans}
\babelfont{tt}{DejaVu Sans Mono}

% =========================================================
% Математика
% =========================================================
\usepackage{amsmath,amssymb}
\usepackage{bm}

% =========================================================
% Графика и таблицы
% =========================================================
\usepackage{graphicx}
\usepackage{booktabs}
\usepackage{array}

% =========================================================
% TikZ
% =========================================================
\usepackage{tikz}
\usetikzlibrary{shapes.geometric, arrows.meta, positioning, calc, fit, backgrounds}

% =========================================================
% Тема
% =========================================================
\usetheme{default}
\usecolortheme{default}

% =========================================================
% Цвета СВФУ
% =========================================================
\definecolor{svfuDark}{RGB}{10,40,90}
\definecolor{svfuAccent}{RGB}{30,144,255}
\definecolor{svfuLight}{RGB}{230,240,255}
\definecolor{sqlLight}{RGB}{245,248,252}

% =========================================================
% Стили TikZ
% =========================================================
\tikzset{
  entity/.style={
    rectangle, draw=svfuDark, very thick, fill=svfuLight,
    minimum width=3.0cm, minimum height=0.9cm,
    font=\bfseries\small, align=center, rounded corners=1pt
  },
  relationship/.style={
    diamond, draw=svfuDark, very thick, fill=svfuAccent!25,
    aspect=2.4, font=\scriptsize, align=center, inner sep=1pt
  },
  module/.style={
    rectangle, draw=svfuDark, very thick, fill=white,
    minimum width=3.1cm, minimum height=0.95cm,
    font=\footnotesize\bfseries, align=center, rounded corners=2pt
  },
  actor/.style={
    rectangle, draw=svfuDark, thick, fill=svfuLight,
    minimum width=2.6cm, minimum height=0.85cm,
    font=\small\bfseries, align=center, rounded corners=4pt
  },
  cmd/.style={
    rectangle, draw=svfuDark, thick, fill=sqlLight,
    inner sep=5pt, font=\ttfamily\small, align=left
  },
  phase/.style={
    rectangle, draw=svfuDark, thick, fill=svfuLight,
    rounded corners=2pt, minimum width=2.5cm, minimum height=0.8cm,
    font=\small\bfseries, align=center
  }
}

% =========================================================
% Основные цвета Beamer
% =========================================================
\setbeamercolor{normal text}{fg=black,bg=white}
\setbeamercolor{structure}{fg=svfuDark}
\setbeamercolor{alerted text}{fg=svfuDark}
\setbeamercolor{frametitle}{bg=svfuDark,fg=white}
\setbeamercolor{title}{fg=white}
\setbeamercolor{author}{fg=white}
\setbeamercolor{date}{fg=white}
\setbeamercolor{accent line}{bg=svfuAccent}
\setbeamercolor{footline left}{bg=svfuDark,fg=white}
\setbeamercolor{footline center}{bg=svfuDark!85,fg=white!80}
\setbeamercolor{footline right}{bg=svfuDark!70,fg=white}
\setbeamercolor{block title}{bg=svfuDark,fg=white}
\setbeamercolor{block body}{bg=svfuLight,fg=black}

% =========================================================
% Шрифты и шаблоны
% =========================================================
\setbeamerfont{frametitle}{size=\large,series=\bfseries}
\setbeamerfont{title}{size=\Large,series=\bfseries}
\setbeamertemplate{navigation symbols}{}
\setbeamertemplate{headline}{}

\setbeamertemplate{frametitle}{%
  \vspace*{-1pt}%
  \begin{beamercolorbox}[wd=\paperwidth,ht=3.2ex,dp=1.2ex,leftskip=10pt,sep=0pt]{frametitle}%
    \usebeamerfont{frametitle}\insertframetitle
  \end{beamercolorbox}%
}

\setbeamertemplate{footline}{%
  \begin{beamercolorbox}[wd=\paperwidth,ht=0.4pt,dp=0pt]{accent line}\end{beamercolorbox}%
  \hbox to \paperwidth{%
    \begin{beamercolorbox}[wd=0.305\paperwidth,ht=3ex,dp=1.5ex,left,sep=0pt,leftskip=6ex]{footline left}
      \textbf{СВФУ}
    \end{beamercolorbox}%
    \begin{beamercolorbox}[wd=0.295\paperwidth,ht=3ex,dp=1.5ex,center,sep=0pt]{footline center}\end{beamercolorbox}%
    \begin{beamercolorbox}[wd=0.4\paperwidth,ht=3ex,dp=1.5ex,right,sep=0pt]{footline right}
      \raggedleft\today\hspace{5ex}\insertframenumber\,/\,\inserttotalframenumber\hspace{6ex}
    \end{beamercolorbox}%
  }%
}

\setbeamertemplate{itemize item}{\color{svfuAccent}\small$\blacktriangleright$}
\setbeamertemplate{itemize subitem}{\color{svfuAccent}\scriptsize$\blacktriangleright$}
\setlength{\leftmargini}{1.5em}

\newcommand{\codebox}[1]{%
  \begin{center}
  \fcolorbox{svfuDark}{sqlLight}{\begin{minipage}{0.90\textwidth}\vspace{2pt}\small\ttfamily #1\vspace{2pt}\end{minipage}}
  \end{center}}

\begin{document}

% =========================================================
% ТИТУЛ
% =========================================================
\begin{frame}[plain]
  \vspace*{1cm}
  \begin{center}
    \colorbox{svfuDark}{%
      \parbox[c][2.8cm][c]{0.86\paperwidth}{%
        \centering
        {\LARGE\color{white}Лекция 2. Системы управления базами данных и язык SQL}
      }%
    }
  \end{center}
  \vfill
  \centering
  {\normalsize Дисциплина: Методы и средства проектирования баз данных}\\[4pt]
  {\small Составитель: Монастырев Егорий Егорьевич, М-ВТ-25}\\[4pt]
  {\small Якутск, 2026}
  \vfill
\end{frame}

% =========================================================
% 1. ЦЕЛИ
% =========================================================
\begin{frame}{Что должно стать понятным после лекции}
\begin{itemize}
  \item чем СУБД отличается от самой базы данных и какую роль она выполняет;
  \item из каких основных компонентов состоит СУБД и как проходит запрос;
  \item как клиентское приложение взаимодействует с сервером базы данных;
  \item что такое SQL и как его команды группируются по назначению;
  \item чем отличаются DDL, DML, DCL и TCL;
  \item как выглядят базовые операции PostgreSQL, которые будут использованы в лабораторной работе.
\end{itemize}

\vspace{0.25cm}
\begin{alertblock}{Логика лекции}
От устройства СУБД \;$\rightarrow$\; к языку SQL \;$\rightarrow$\; к конкретным командам.
\end{alertblock}
\end{frame}

% =========================================================
% 2. ПЕРЕХОД ОТ ЛЕКЦИИ 1
% =========================================================
\begin{frame}{Что мы уже знаем из лекции 1}
\small
В лекции 1 мы рассмотрели предметную область, данные и основные модели их представления,
ER-модель и переход от предметной области к структуре базы данных.

\vspace{0.35cm}
\begin{center}
\begin{tikzpicture}[node distance=0.35cm and 0.7cm]
  \node[phase] (domain) {Предметная\\область};
  \node[phase,right=of domain] (model) {Модель\\данных};
  \node[phase,right=of model] (db) {База\\данных};
  \node[phase,right=of db,minimum width=3.1cm] (dbms) {СУБД};
  \draw[-{Latex[length=2.2mm]},thick,svfuAccent] (domain)--(model);
  \draw[-{Latex[length=2.2mm]},thick,svfuAccent] (model)--(db);
  \draw[-{Latex[length=2.2mm]},thick,svfuAccent] (db)--(dbms);
\end{tikzpicture}
\end{center}

\vspace{0.3cm}
\begin{alertblock}{Новая точка фокуса}
Не повторяем построение ER-модели и архитектуру ANSI/SPARC. Теперь рассматриваем,
\textbf{как СУБД реально обслуживает операции над уже спроектированной БД}.
\end{alertblock}
\end{frame}

% =========================================================
% 3. СУБД
% =========================================================
\begin{frame}{СУБД: не просто «хранилище таблиц»}
\begin{block}{Определение}
\textbf{СУБД} — программная система, предназначенная для определения структуры базы данных,
хранения и обработки данных, управления доступом к ним, поддержания целостности,
транзакционной обработки и восстановления после сбоев.
\end{block}

\vspace{0.35cm}
\begin{itemize}
  \item приложение формулирует операцию;
  \item СУБД проверяет права и ограничения;
  \item СУБД выбирает способ выполнения операции;
  \item СУБД организует чтение или изменение данных;
  \item результат возвращается клиенту.
\end{itemize}
\end{frame}

\begin{frame}{Основные функции СУБД}
\small
\begin{table}
\centering
\begin{tabular}{p{0.31\textwidth}|p{0.57\textwidth}}
\toprule
\textbf{Функция} & \textbf{Задача}\\
\midrule
Хранение & размещение и чтение данных на внешней памяти\\
Обработка запросов & разбор, построение и выполнение плана запроса\\
Целостность & контроль ограничений и допустимых состояний данных\\
Транзакции & согласованное выполнение группы операций\\
Конкурентный доступ & корректная работа нескольких пользователей\\
Безопасность & аутентификация, роли и права доступа\\
Журналирование и восстановление & сохранение информации, необходимой для восстановления после сбоев\\
\bottomrule
\end{tabular}
\end{table}
\end{frame}

% =========================================================
% 4. АРХИТЕКТУРА
% =========================================================
\begin{frame}{Из каких частей состоит СУБД}
\begin{center}
\begin{tikzpicture}[node distance=0.55cm and 1.0cm]
  \node[module,minimum width=3.4cm] (parser) {Обработчик\\запросов};
  \node[module,below=of parser] (tx) {Менеджер\\транзакций};
  \node[module,below=of tx] (storage) {Менеджер\\хранения};
  \node[module,right=1.0cm of parser] (catalog) {Системный каталог\\и метаданные};
  \node[module,right=1.0cm of tx] (buffer) {Буферный\\менеджер};
  \node[module,right=1.0cm of storage] (recovery) {Журналирование\\и восстановление};
  \draw[thick,svfuAccent] (parser)--(catalog);
  \draw[thick,svfuAccent] (parser)--(tx);
  \draw[thick,svfuAccent] (tx)--(buffer);
  \draw[thick,svfuAccent] (storage)--(buffer);
  \draw[thick,svfuAccent] (tx)--(recovery);
\end{tikzpicture}
\end{center}
\vspace{0.2cm}
\begin{alertblock}{Не нужно запоминать реализацию конкретного продукта}
Архитектура PostgreSQL, MySQL, Oracle и других СУБД различается деталями,
но набор решаемых задач имеет одну общую инженерную основу.
\end{alertblock}
\end{frame}

\begin{frame}{Путь SQL-запроса через СУБД}
\begin{center}
\begin{tikzpicture}[node distance=0.55cm and 0.55cm]
  \node[actor] (app) {Клиент\texttt{SELECT ...}};
  \node[module,right=of app] (parser) {Разбор\\запроса};
  \node[module,right=of parser] (plan) {Планирование\\и оптимизация};
  \node[module,right=of plan] (exec) {Выполнение};
  \node[entity,right=of exec,minimum width=2.3cm] (data) {Данные};
  \draw[-{Latex[length=2.3mm]},thick,svfuAccent] (app)--(parser);
  \draw[-{Latex[length=2.3mm]},thick,svfuAccent] (parser)--(plan);
  \draw[-{Latex[length=2.3mm]},thick,svfuAccent] (plan)--(exec);
  \draw[-{Latex[length=2.3mm]},thick,svfuAccent] (exec)--(data);
\end{tikzpicture}
\end{center}

\vspace{0.35cm}
\begin{itemize}
  \item разбор: синтаксис и смысловые элементы;
  \item планирование: какой способ выполнения выбрать;
  \item выполнение: чтение или изменение данных;
  \item возврат результата клиенту.
\end{itemize}
\end{frame}

\begin{frame}{Клиент и сервер базы данных}
\begin{center}
\begin{tikzpicture}[node distance=0.45cm and 1.0cm]
  \node[actor] (gui) {Приложение / CLI-клиент};
  \node[module,right=of gui,minimum width=4.0cm] (server) {Сервер СУБД\footnotesize PostgreSQL};
  \node[entity,right=of server,minimum width=3.0cm] (db) {База данных};
  \draw[<->,thick,svfuAccent] (gui)--(server);
  \draw[<->,thick,svfuDark] (server)--(db);
\end{tikzpicture}
\end{center}

\vspace{0.35cm}
\begin{block}{Ключевая идея}
Клиент не работает непосредственно с файлами базы данных. Он отправляет запросы
серверу СУБД, а сервер отвечает за выполнение операций и защиту данных.
\end{block}
\end{frame}

\begin{frame}{Кто взаимодействует с СУБД}
\small
\begin{table}
\centering
\begin{tabular}{p{0.32\textwidth}|p{0.56\textwidth}}
\toprule
\textbf{Участник} & \textbf{Типичные задачи}\\
\midrule
Администратор БД & создание ролей, настройка доступа, резервное копирование, контроль эксплуатации\\
Разработчик & проектирование схемы, SQL, транзакции, работа с API СУБД\\
Аналитик / специалист & формирование запросов, выборка и анализ данных\\
Конечный пользователь & работа с данными через прикладное приложение\\
\bottomrule
\end{tabular}
\end{table}
\end{frame}

% =========================================================
% 5. SQL
% =========================================================
\begin{frame}{SQL как язык взаимодействия с реляционной СУБД}
\begin{block}{SQL}
\textbf{SQL (Structured Query Language)} — стандартизованный язык работы
с реляционными базами данных, включающий средства определения структуры,
манипулирования данными и управления доступом и транзакциями.
\end{block}

\vspace{0.3cm}
\begin{itemize}
  \item SQL в основном декларативен: описывается требуемый результат или операция;
  \item конкретный план выполнения выбирает СУБД;
  \item разные СУБД поддерживают SQL с диалектными и функциональными отличиями.
\end{itemize}
\end{frame}

\begin{frame}{Подъязыки SQL: карта лекции}
\begin{center}
\begin{tikzpicture}[node distance=0.4cm and 0.55cm]
  \node[phase,minimum width=2.7cm] (ddl) {DDL\\\footnotesize структура};
  \node[phase,right=of ddl,minimum width=2.7cm] (dml) {DML\\\footnotesize данные};
  \node[phase,right=of dml,minimum width=2.7cm] (dcl) {DCL\\\footnotesize доступ};
  \node[phase,right=of dcl,minimum width=2.7cm] (tcl) {TCL\\\footnotesize транзакции};
\end{tikzpicture}
\end{center}

\vspace{0.45cm}
\small
\begin{tabular}{p{0.14\textwidth}|p{0.38\textwidth}|p{0.34\textwidth}}
\toprule
\textbf{DDL} & определение объектов БД & \texttt{CREATE}, \texttt{ALTER}, \texttt{DROP}\\
\textbf{DML} & изменение и выборка данных & \texttt{SELECT}, \texttt{INSERT}, \texttt{UPDATE}, \texttt{DELETE}\\
\textbf{DCL} & управление правами & \texttt{GRANT}, \texttt{REVOKE}\\
\textbf{TCL} & управление транзакциями & \texttt{COMMIT}, \texttt{ROLLBACK}, \texttt{SAVEPOINT}\\
\bottomrule
\end{tabular}

\vspace{0.2cm}
\footnotesize
Примечание: классификация команд в учебной литературе может отличаться; например,
\texttt{SELECT} иногда выделяют отдельно как DQL.
\end{frame}

% =========================================================
% DDL
% =========================================================
\begin{frame}{DDL: Data Definition Language}
\begin{block}{Назначение}
DDL используется для \textbf{определения и изменения объектов базы данных}:
схем, таблиц, ограничений, индексов, представлений и других объектов.
\end{block}

\vspace{0.3cm}
\begin{center}
\begin{tikzpicture}[node distance=0.55cm and 0.55cm]
  \node[cmd] (c) {CREATE};
  \node[cmd,right=of c] (a) {ALTER};
  \node[cmd,right=of a] (d) {DROP};
  \draw[-{Latex[length=2mm]},thick,svfuAccent] (c)--(a);
  \draw[-{Latex[length=2mm]},thick,svfuAccent] (a)--(d);
\end{tikzpicture}
\end{center}

\begin{alertblock}{Практический смысл}
DDL отвечает на вопрос: \textbf{какие объекты должны существовать в БД и как они устроены?}
\end{alertblock}
\end{frame}

% =========================================================
% ТИПЫ ДАННЫХ
% =========================================================
\begin{frame}{Типы данных в PostgreSQL}
\begin{block}{Зачем нужен тип данных?}
Тип данных определяет, какие значения может хранить столбец,
какие операции над ними допустимы и как СУБД должна их интерпретировать.
\end{block}

\vspace{0.3cm}
\small
\begin{table}
\centering
\begin{tabular}{p{0.28\textwidth}|p{0.25\textwidth}|p{0.34\textwidth}}
\toprule
\textbf{Группа} & \textbf{Примеры} & \textbf{Для чего}\\
\midrule
Числовые & \texttt{INTEGER}, \texttt{BIGINT}, \texttt{NUMERIC} & идентификаторы, счётчики, деньги\\
Строковые & \texttt{TEXT}, \texttt{VARCHAR(n)} & имена, коды, описания\\
Дата и время & \texttt{DATE}, \texttt{TIMESTAMP} & даты, моменты времени\\
Логические & \texttt{BOOLEAN} & да/нет, включено/выключено\\
Специальные & \texttt{UUID}, \texttt{JSONB} & идентификаторы и полуструктурированные данные\\
\bottomrule
\end{tabular}
\end{table}

\vspace{0.2cm}
\alert{Выбор типа — часть проектирования схемы, а не декоративное свойство столбца.}
\end{frame}

\begin{frame}{Числовые типы данных}
\small
\begin{table}
\centering
\begin{tabular}{p{0.25\textwidth}|p{0.29\textwidth}|p{0.34\textwidth}}
\toprule
\textbf{Тип} & \textbf{Характеристика} & \textbf{Пример}\\
\midrule
\texttt{SMALLINT} & малый целый тип & номер курса, небольшой счётчик\\
\texttt{INTEGER} & основной целый тип & идентификатор, количество\\
\texttt{BIGINT} & большой целый тип & большие счётчики и идентификаторы\\
\texttt{NUMERIC(p,s)} & точная десятичная арифметика & цена, сумма, финансовые значения\\
\texttt{REAL} / \texttt{DOUBLE PRECISION} & числа с плавающей точкой & измерения, научные вычисления\\
\bottomrule
\end{tabular}
\end{table}

\vspace{0.25cm}
\begin{alertblock}{Практическое правило}
Для денег и других величин, где важна точность десятичного представления,
обычно выбирают \texttt{NUMERIC}, а не тип с плавающей точкой.
\end{alertblock}
\end{frame}

\begin{frame}{Строковые, логические и временные типы}
\small
\begin{table}
\centering
\begin{tabular}{p{0.29\textwidth}|p{0.30\textwidth}|p{0.29\textwidth}}
\toprule
\textbf{Тип} & \textbf{Пример значения} & \textbf{Пример применения}\\
\midrule
\texttt{TEXT} & \texttt{'Иванов Иван'} & имя, описание, текст\\
\texttt{VARCHAR(n)} & \texttt{'ПМИ-25'} & строка с ограниченной длиной\\
\texttt{BOOLEAN} & \texttt{TRUE} / \texttt{FALSE} & активен ли пользователь\\
\texttt{DATE} & \texttt{2026-09-07} & дата рождения, дата события\\
\texttt{TIMESTAMP} & дата + время & момент регистрации события\\
\texttt{TIMESTAMPTZ} & дата + время с часовым поясом & события в распределённых системах\\
\bottomrule
\end{tabular}
\end{table}

\vspace{0.25cm}
\begin{block}{Важно}
\texttt{DATE} хранит дату, а \texttt{TIMESTAMP/TIMESTAMPTZ} — момент или дату с временной составляющей.
Выбирать тип нужно исходя из смысла данных.
\end{block}
\end{frame}

\begin{frame}{Специальные типы: UUID и JSONB}
\small
\begin{block}{UUID}
\texttt{UUID} — 128-битный идентификатор, представляемый шестнадцатеричной строкой.
Удобен, когда идентификаторы должны быть уникальны независимо от локальной последовательности чисел.
\end{block}

\vspace{0.2cm}
\codebox{CREATE TABLE device (\\
\hspace*{2em}id UUID PRIMARY KEY,\\
\hspace*{2em}config JSONB\\
);}

\vspace{0.15cm}
\begin{block}{JSONB}
\texttt{JSONB} позволяет хранить структурированные документы в бинарно обработанном
внутреннем представлении PostgreSQL. Это не отменяет реляционное проектирование:
таблицы остаются основным средством моделирования строгих связей и ограничений.
\end{block}
\end{frame}

\begin{frame}{Как выбирать тип данных}
\begin{enumerate}
  \item Определить \textbf{смысл} значения, а не только его внешний вид.
  \item Выбрать тип, который естественно выражает этот смысл.
  \item Не хранить числа и даты как \texttt{TEXT}, если над ними нужны соответствующие операции.
  \item Для денежных и точных десятичных значений использовать \texttt{NUMERIC}.
  \item Ограничения \texttt{NOT NULL}, \texttt{CHECK}, \texttt{UNIQUE} и ключи задавать отдельно — тип данных не заменяет ограничения целостности.
\end{enumerate}

\vspace{0.25cm}
\begin{alertblock}{Пример}
\texttt{age INTEGER} лучше, чем \texttt{age TEXT};
\texttt{birth\_date DATE} лучше, чем \texttt{birth\_date TEXT}.
СУБД сможет корректно сравнивать, сортировать и проверять такие значения.
\end{alertblock}
\end{frame}

\begin{frame}{Тип данных и значение по умолчанию}
\small
Тип определяет множество допустимых значений,
а \texttt{DEFAULT} определяет значение, которое СУБД может подставить,
если при вставке оно не указано.

\vspace{0.25cm}
\codebox{CREATE TABLE student (\\
\hspace*{2em}id INTEGER GENERATED ALWAYS AS IDENTITY PRIMARY KEY,\\
\hspace*{2em}name TEXT NOT NULL,\\
\hspace*{2em}course INTEGER DEFAULT 1,\\
\hspace*{2em}active BOOLEAN DEFAULT TRUE\\
);}

\vspace{0.2cm}
\begin{itemize}
  \item \texttt{INTEGER}, \texttt{TEXT}, \texttt{BOOLEAN} — типы данных;
  \item \texttt{DEFAULT} — значение по умолчанию;
  \item \texttt{NOT NULL} — запрет отсутствующего значения;
  \item \texttt{PRIMARY KEY} — ограничение уникальной идентификации.
\end{itemize}
\end{frame}

\begin{frame}{NULL: отсутствие значения — не тип данных}
\small
\begin{block}{Что такое \texttt{NULL}?}
\texttt{NULL} обозначает отсутствие известного или заданного значения.
Это не строка \texttt{'NULL'}, не число 0 и не пустая строка \texttt{''}.
\end{block}

\vspace{0.25cm}
\begin{table}
\centering
\begin{tabular}{p{0.23\textwidth}|p{0.27\textwidth}|p{0.36\textwidth}}
\toprule
\textbf{Значение} & \textbf{Пример} & \textbf{Смысл}\\
\midrule
Число & \texttt{0} & известное числовое значение\\
Пустая строка & \texttt{''} & известная строка длины 0\\
\texttt{NULL} & \texttt{NULL} & значение отсутствует / неизвестно\\
\bottomrule
\end{tabular}
\end{table}

\vspace{0.2cm}
\alert{Проверка отсутствия значения выполняется через \texttt{IS NULL} / \texttt{IS NOT NULL}, а не через \texttt{= NULL}.}
\end{frame}

\begin{frame}{Идентификаторы: SERIAL и IDENTITY}
\small
В учебных примерах часто встречается:
\codebox{CREATE TABLE student (\\
\hspace*{2em}id SERIAL PRIMARY KEY,\\
\hspace*{2em}name TEXT NOT NULL\\
);}

\vspace{0.2cm}
\begin{block}{Современный вариант PostgreSQL}
Для автоматически генерируемых идентификаторов рекомендуется рассматривать стандартный механизм
\texttt{GENERATED ... AS IDENTITY}:
\end{block}

\codebox{CREATE TABLE student (\
\hspace*{2em}id INTEGER GENERATED ALWAYS AS IDENTITY PRIMARY KEY,\
\hspace*{2em}name TEXT NOT NULL\
);}

\vspace{0.15cm}
\footnotesize
\texttt{SERIAL} — удобный исторический механизм на основе последовательности;
\texttt{IDENTITY} — более явный современный синтаксис определения автоматически генерируемого столбца.
\end{frame}

\begin{frame}{DDL: создание таблицы}
\codebox{CREATE TABLE student (\\
\hspace*{2em}id INTEGER PRIMARY KEY,\\
\hspace*{2em}name TEXT NOT NULL,\\
\hspace*{2em}group\_id INTEGER\\
);}

\vspace{0.25cm}
\begin{itemize}
  \item имя таблицы задаёт объект схемы;
  \item столбцы определяют структуру строки;
  \item тип данных ограничивает множество допустимых значений;
  \item ограничения задают правила целостности.
\end{itemize}
\end{frame}

\begin{frame}{DDL: типы данных и ограничения}
\small
\begin{table}
\centering
\begin{tabular}{p{0.27\textwidth}|p{0.25\textwidth}|p{0.36\textwidth}}
\toprule
\textbf{Элемент} & \textbf{Пример} & \textbf{Идея}\\
\midrule
Тип & \texttt{INTEGER}, \texttt{TEXT}, \texttt{DATE} & какие значения допустимы\\
\texttt{PRIMARY KEY} & \texttt{id} & уникальный идентификатор строки\\
\texttt{NOT NULL} & \texttt{name} & значение обязательно\\
\texttt{UNIQUE} & \texttt{email} & значения не должны повторяться\\
\texttt{CHECK} & \texttt{score BETWEEN 2 AND 5} & произвольное условие\\
\texttt{FOREIGN KEY} & \texttt{group\_id} & ссылочная целостность\\
\bottomrule
\end{tabular}
\end{table}

\vspace{0.2cm}
\alert{DDL определяет правила, а не сами экземпляры данных.}
\end{frame}

\begin{frame}{DDL: изменение структуры с помощью ALTER}
\codebox{ALTER TABLE student ADD COLUMN email TEXT;}
\vspace{0.15cm}
\codebox{ALTER TABLE student ADD CONSTRAINT student\_email\_uq UNIQUE (email);}

\vspace{0.25cm}
\begin{itemize}
  \item \texttt{ADD COLUMN} — добавить столбец;
  \item \texttt{ALTER COLUMN} — изменить свойства столбца;
  \item \texttt{ADD CONSTRAINT} — добавить ограничение;
  \item состав конкретных возможностей зависит от СУБД и её версии.
\end{itemize}
\end{frame}

\begin{frame}{DDL: удаление объектов и осторожность}
\begin{center}
\begin{tikzpicture}[node distance=0.6cm]
  \node[cmd] (drop) {DROP TABLE student;};
  \node[cmd,below=of drop] (alter) {ALTER TABLE student ...;};
  \node[cmd,below=of alter] (create) {CREATE TABLE student (...);};
  \draw[-{Latex[length=2mm]},thick,svfuAccent] (drop)--(alter);
  \draw[-{Latex[length=2mm]},thick,svfuAccent] (alter)--(create);
\end{tikzpicture}
\end{center}

\vspace{0.25cm}
\begin{alertblock}{Операции DDL требуют внимания}
Удаление или изменение структуры может сделать существующие запросы и приложения
несовместимыми со схемой БД. Поэтому DDL — часть проектирования и сопровождения схемы.
\end{alertblock}
\end{frame}

\begin{frame}{DDL: небольшой пример схемы «Университет»}
\codebox{CREATE TABLE student (\\
\hspace*{2em}id INTEGER PRIMARY KEY,\\
\hspace*{2em}name TEXT NOT NULL,\\
\hspace*{2em}group\_id INTEGER REFERENCES student\_group(id)\\
);}

\vspace{0.2cm}
\begin{block}{Что делает этот код?}
Он создаёт таблицу и одновременно фиксирует часть логической модели:
у студента есть идентификатор, имя и ссылка на группу.
\end{block}

\footnotesize
Полная схема должна быть создана в согласованном порядке, учитывая зависимости внешних ключей.
\end{frame}

% =========================================================
% DML
% =========================================================
\begin{frame}{DML: Data Manipulation Language}
\begin{block}{Назначение}
DML используется для \textbf{работы с данными внутри уже созданных объектов БД}.
\end{block}

\vspace{0.3cm}
\begin{center}
\begin{tikzpicture}[node distance=0.55cm and 0.55cm]
  \node[cmd] (i) {INSERT};
  \node[cmd,right=of i] (s) {SELECT};
  \node[cmd,right=of s] (u) {UPDATE};
  \node[cmd,right=of u] (d) {DELETE};
\end{tikzpicture}
\end{center}

\vspace{0.35cm}
\begin{alertblock}{Важное различие}
DDL меняет \textbf{структуру объектов}; DML работает с \textbf{экземплярами данных}.
\end{alertblock}
\end{frame}

\begin{frame}{DML: добавление данных — INSERT}
\codebox{INSERT INTO student (id, name, group\_id)\\
VALUES (101, 'Иванов Иван', 10);}

\vspace{0.25cm}
\begin{itemize}
  \item явно перечислять столбцы — надёжнее, чем полагаться на порядок всех полей;
  \item СУБД проверяет типы и ограничения;
  \item при нарушении ограничения операция завершается ошибкой.
\end{itemize}
\end{frame}

\begin{frame}{DML: выборка данных — SELECT}
\codebox{SELECT id, name\\
FROM student\\
WHERE group\_id = 10;}

\vspace{0.25cm}
\begin{itemize}
  \item \texttt{SELECT} определяет, какие значения нужно получить;
  \item \texttt{FROM} указывает источник данных;
  \item \texttt{WHERE} задаёт условие отбора строк.
\end{itemize}

\begin{alertblock}{Для лабораторной}
На этом синтаксисе начинается практическая работа с запросами.
\end{alertblock}
\end{frame}

\begin{frame}{DML: условия выборки — WHERE}
\small
\begin{table}
\centering
\begin{tabular}{p{0.24\textwidth}|p{0.55\textwidth}}
\toprule
Оператор & Пример\\
\midrule
Сравнение & \texttt{score >= 4}\\
И & \texttt{group\_id = 10 AND name IS NOT NULL}\\
Или & \texttt{group\_id = 10 OR group\_id = 11}\\
Диапазон & \texttt{score BETWEEN 3 AND 5}\\
Множество & \texttt{group\_id IN (10,11,12)}\\
Шаблон & \texttt{name LIKE 'Ив\%'}\\
Пустое значение & \texttt{email IS NULL}\\
\bottomrule
\end{tabular}
\end{table}

\vspace{0.15cm}
\footnotesize
\textbf{Важно:} \texttt{NULL} не равен обычному значению и проверяется через
\texttt{IS NULL}/\texttt{IS NOT NULL}.
\end{frame}

\begin{frame}{DML: изменение данных — UPDATE}
\codebox{UPDATE student\\
SET name = 'Иванов И. И.'\\
WHERE id = 101;}

\vspace{0.3cm}
\begin{alertblock}{Критически важно}
Запрос \texttt{UPDATE} без корректного условия \texttt{WHERE} может изменить
\textbf{все строки} таблицы.
\end{alertblock}

\vspace{0.15cm}
То же правило касается и других массовых операций над данными.
\end{frame}

\begin{frame}{DML: удаление данных — DELETE}
\codebox{DELETE FROM student\\
WHERE id = 101;}

\vspace{0.3cm}
\begin{itemize}
  \item удаляются строки, удовлетворяющие условию;
  \item ограничения внешних ключей могут запретить удаление;
  \item удаление данных и удаление таблицы — разные операции.
\end{itemize}

\begin{block}{Сравнение}
\texttt{DELETE} — работа с данными.\\
\texttt{DROP TABLE} — удаление объекта схемы.
\end{block}
\end{frame}

\begin{frame}{SELECT, INSERT, UPDATE, DELETE: единый цикл работы}
\begin{center}
\begin{tikzpicture}[node distance=0.45cm and 0.5cm]
  \node[phase] (s) {SELECT\\просмотреть};
  \node[phase,right=of s] (i) {INSERT\\добавить};
  \node[phase,right=of i] (u) {UPDATE\\изменить};
  \node[phase,right=of u] (d) {DELETE\\удалить};
  \draw[-{Latex[length=2mm]},thick,svfuAccent] (s)--(i);
  \draw[-{Latex[length=2mm]},thick,svfuAccent] (i)--(u);
  \draw[-{Latex[length=2mm]},thick,svfuAccent] (u)--(d);
\end{tikzpicture}
\end{center}

\vspace{0.35cm}
\begin{alertblock}{Практическая схема}
Создали структуру через DDL \;$\rightarrow$\; заполнили и изменяем её через DML
\;$\rightarrow$\; доступ ограничиваем через DCL \;$\rightarrow$\; согласованность операций управляем через TCL.
\end{alertblock}
\end{frame}

% =========================================================
% DCL
% =========================================================
\begin{frame}{DCL: Data Control Language}
\begin{block}{Назначение}
DCL используется для управления \textbf{правами доступа} пользователей и ролей
к объектам базы данных.
\end{block}

\vspace{0.3cm}
\begin{center}
\begin{tikzpicture}[node distance=0.75cm and 1.4cm]
  \node[actor] (role) {Роль / пользователь};
  \node[module,right=of role,minimum width=3.6cm] (obj) {Объект БД};
  \node[cmd,below=of obj] (grant) {GRANT / REVOKE};
  \draw[<->,thick,svfuAccent] (role)--(obj);
  \draw[-{Latex[length=2mm]},thick,svfuDark] (grant)--(obj);
\end{tikzpicture}
\end{center}
\end{frame}

\begin{frame}{DCL: GRANT — выдать право}
\codebox{GRANT SELECT ON student TO analyst;}
\vspace{0.25cm}
\codebox{GRANT SELECT, INSERT ON student TO operator;}

\vspace{0.25cm}
\begin{itemize}
  \item право задаётся для конкретной роли или пользователя;
  \item можно разрешать отдельные операции;
  \item принцип наименьших привилегий: выдавать только необходимые права.
\end{itemize}
\end{frame}

\begin{frame}{DCL: REVOKE — забрать право}
\codebox{REVOKE INSERT ON student FROM analyst;}

\vspace{0.3cm}
\begin{itemize}
  \item \texttt{GRANT} добавляет разрешение;
  \item \texttt{REVOKE} отзывает разрешение;
  \item система прав может быть многоуровневой и включать наследование ролей.
\end{itemize}

\begin{alertblock}{Зачем это нужно}
Даже корректно спроектированная БД небезопасна, если любой пользователь может
изменять любые данные.
\end{alertblock}
\end{frame}

\begin{frame}{DCL: пример разделения доступа}
\small
\begin{table}
\centering
\begin{tabular}{p{0.30\textwidth}|p{0.20\textwidth}|p{0.38\textwidth}}
\toprule
\textbf{Роль} & \textbf{Право} & \textbf{Назначение}\\
\midrule
\texttt{analyst} & SELECT & анализ данных без изменения\\
\texttt{operator} & SELECT, INSERT, UPDATE & операционная работа\\
\texttt{auditor} & SELECT & контроль и аудит\\
\texttt{admin} & расширенный доступ & администрирование\\
\bottomrule
\end{tabular}
\end{table}

\vspace{0.3cm}
\begin{block}{Принцип}
Права должны отражать реальную роль человека или сервиса, а не выдаваться
«на всякий случай».
\end{block}
\end{frame}

% =========================================================
% TCL
% =========================================================
\begin{frame}{TCL: управление транзакциями}
\begin{block}{Транзакция}
\textbf{Транзакция} — логически связанная последовательность операций,
которую СУБД рассматривает как единое целое для фиксации или отмены результата.
\end{block}

\vspace{0.25cm}
\begin{center}
\begin{tikzpicture}[node distance=0.5cm and 0.65cm]
  \node[phase] (begin) {BEGIN};
  \node[phase,right=of begin] (ops) {Операции};
  \node[phase,right=of ops] (commit) {COMMIT};
  \node[phase,below=0.55cm of ops] (rollback) {ROLLBACK};
  \draw[-{Latex[length=2mm]},thick,svfuAccent] (begin)--(ops);
  \draw[-{Latex[length=2mm]},thick,svfuAccent] (ops)--(commit);
  \draw[-{Latex[length=2mm]},thick,svfuDark] (ops)--(rollback);
\end{tikzpicture}
\end{center}
\end{frame}

\begin{frame}{TCL: COMMIT и ROLLBACK}
\codebox{BEGIN;\\
UPDATE student SET group\_id = 11 WHERE id = 101;\\
COMMIT;}

\vspace{0.2cm}
\codebox{BEGIN;\\
UPDATE student SET group\_id = 99 WHERE id = 101;\\
ROLLBACK;}

\vspace{0.2cm}
\begin{itemize}
  \item \texttt{COMMIT} фиксирует изменения транзакции;
  \item \texttt{ROLLBACK} отменяет незафиксированные изменения транзакции.
\end{itemize}
\end{frame}

\begin{frame}{TCL: SAVEPOINT}
\codebox{BEGIN;\\
INSERT INTO student VALUES (102, 'Петров П. П.', 10);\\
SAVEPOINT before;\\
UPDATE student SET group\_id = 11 WHERE id = 102;\\
ROLLBACK TO SAVEPOINT before;\\
COMMIT;}

\vspace{0.25cm}
\begin{block}{Идея}
\texttt{SAVEPOINT} создаёт промежуточную точку, к которой можно вернуться,
не отменяя всю транзакцию целиком.
\end{block}
\end{frame}

\begin{frame}{Транзакция на примере перевода}
\begin{center}
\begin{tikzpicture}[node distance=0.55cm and 0.65cm]
  \node[actor] (a) {Счёт A\\1000};
  \node[module,right=of a] (s1) {списать 200};
  \node[module,right=of s1] (s2) {зачислить 200};
  \node[actor,right=of s2] (b) {Счёт B\\500};
  \draw[-{Latex[length=2mm]},thick,svfuAccent] (a)--(s1);
  \draw[-{Latex[length=2mm]},thick,svfuAccent] (s1)--(s2);
  \draw[-{Latex[length=2mm]},thick,svfuAccent] (s2)--(b);
\end{tikzpicture}
\end{center}

\vspace{0.25cm}
\begin{alertblock}{Почему нужна транзакция?}
Если первая операция выполнена, а вторая завершилась ошибкой, состояние «деньги списали,
но не зачислили» недопустимо. Группа операций должна быть согласована как единое целое.
\end{alertblock}
\end{frame}

\begin{frame}{ACID: почему транзакции являются основой надёжности}
\small
\begin{table}
\centering
\begin{tabular}{p{0.23\textwidth}|p{0.62\textwidth}}
\toprule
\textbf{Свойство} & \textbf{Смысл}\\
\midrule
Atomicity & транзакция выполняется целиком или не даёт своего результата\\
Consistency & после успешной транзакции должны сохраняться заданные правила целостности\\
Isolation & параллельные транзакции должны взаимодействовать контролируемым образом\\
Durability & зафиксированный результат должен сохраняться при штатных сбоях\\
\bottomrule
\end{tabular}
\end{table}

\vspace{0.2cm}
\footnotesize
Детали уровней изоляции, блокировок и многоверсионного контроля — отдельная тема;
здесь фиксируется только логическая основа.
\end{frame}

% =========================================================
% СВОДКА
% =========================================================
\begin{frame}{DDL, DML, DCL и TCL: не путать назначение}
\small
\begin{table}
\centering
\begin{tabular}{p{0.15\textwidth}|p{0.30\textwidth}|p{0.42\textwidth}}
\toprule
\textbf{Группа} & \textbf{Вопрос} & \textbf{Примеры}\\
\midrule
DDL & Что существует в БД? & \texttt{CREATE}, \texttt{ALTER}, \texttt{DROP}\\
DML & Какие данные храним и изменяем? & \texttt{SELECT}, \texttt{INSERT}, \texttt{UPDATE}, \texttt{DELETE}\\
DCL & Кто что может делать? & \texttt{GRANT}, \texttt{REVOKE}\\
TCL & Как зафиксировать или отменить группу операций? & \texttt{COMMIT}, \texttt{ROLLBACK}, \texttt{SAVEPOINT}\\
\bottomrule
\end{tabular}
\end{table}

\vspace{0.35cm}
\begin{center}
\Large
\textbf{Структура} $\rightarrow$ \textbf{Данные} $\rightarrow$ \textbf{Доступ} $\rightarrow$ \textbf{Согласованность}
\end{center}
\end{frame}

\begin{frame}{Минимальный сценарий работы с PostgreSQL}
\begin{center}
\begin{tikzpicture}[node distance=0.45cm and 0.35cm]
  \node[phase] (ddl) {1. DDL\\создать схему};
  \node[phase,right=of ddl] (dml) {2. DML\\записать данные};
  \node[phase,right=of dml] (dcl) {3. DCL\\выдать права};
  \node[phase,right=of dcl] (tcl) {4. TCL\\управлять операциями};
\end{tikzpicture}
\end{center}

\vspace{0.45cm}
\begin{block}{Пример рабочего цикла}
Создать таблицы \;$\rightarrow$\; вставить данные \;$\rightarrow$\; выполнить SELECT
\;$\rightarrow$\; изменить данные в транзакции \;$\rightarrow$\; зафиксировать результат.
\end{block}

\begin{alertblock}{Следующий шаг}
В лабораторной работе этот цикл превращается в практику: студенты начинают писать
и выполнять реальные SQL-запросы к PostgreSQL.
\end{alertblock}
\end{frame}

\begin{frame}{Типичные ошибки начинающего разработчика}
\begin{itemize}
  \item путать таблицу с базой данных, а SQL — с самой СУБД;
  \item менять структуру таблиц, не учитывая существующие зависимости;
  \item выполнять \texttt{UPDATE} или \texttt{DELETE} без корректного \texttt{WHERE};
  \item выдавать роли избыточные права;
  \item выполнять несколько связанных изменений без транзакции;
  \item считать, что успешный запуск команды автоматически означает корректность результата.
\end{itemize}
\end{frame}

\begin{frame}{Что важно запомнить}
\begin{enumerate}
  \item СУБД управляет не только хранением, но и запросами, целостностью, доступом и транзакциями.
  \item SQL является основным языком взаимодействия с реляционной СУБД.
  \item DDL описывает структуру объектов базы данных.
  \item DML выполняет операции над данными.
  \item DCL управляет правами доступа.
  \item TCL управляет фиксацией и отменой транзакций.
  \item Одна прикладная задача часто требует совместного использования всех этих средств.
\end{enumerate}
\end{frame}

\begin{frame}{Контрольные вопросы}
\small
\begin{enumerate}
  \item Какую основную проблему решает СУБД, а не просто набор файлов?
  \item Какие основные компоненты СУБД участвуют в обработке запроса?
  \item Как работает модель «клиент — сервер СУБД»?
  \item Что такое SQL и почему его называют декларативным?
  \item Чем отличается DDL от DML?
  \item В чём разница между \texttt{DELETE} и \texttt{DROP TABLE}?
  \item Для чего используются \texttt{GRANT} и \texttt{REVOKE}?
  \item Что такое транзакция и зачем нужны \texttt{COMMIT} и \texttt{ROLLBACK}?
  \item Для чего нужен \texttt{SAVEPOINT}?
  \item Почему отсутствие корректного \texttt{WHERE} опасно в \texttt{UPDATE} и \texttt{DELETE}? 
\end{enumerate}
\end{frame}

\begin{frame}{Подготовка к лабораторной работе}
\begin{block}{Перед началом работы}
Необходимо понимать четыре группы операций:
\end{block}

\begin{center}
\begin{tabular}{p{0.17\textwidth}|p{0.62\textwidth}}
\toprule
DDL & создать таблицы и ограничения\\
DML & вставить, выбрать, изменить и удалить строки\\
DCL & назначить права ролям при необходимости\\
TCL & объединить связанные изменения в транзакцию\\
\bottomrule
\end{tabular}
\end{center}

\vspace{0.3cm}
\textbf{Лабораторная работа:} практическое выполнение SQL-команд в Microsoft Access 2016.
\end{frame}

\begin{frame}[plain,noframenumbering]
\thispagestyle{empty}
\vspace*{\fill}
\begin{center}
{\LARGE\bfseries Спасибо за внимание!}
\end{center}
\vspace*{\fill}
\end{frame}

\end{document}
